% authority: science_draft
\documentclass[11pt]{article}
\usepackage[T1]{fontenc}
\usepackage[utf8]{inputenc}
\usepackage{lmodern}
\usepackage[margin=1in]{geometry}
\usepackage{amsmath,amssymb,amsthm}
\usepackage{booktabs,longtable,array,enumitem}
\newtheorem{definition}{Definition}
\newtheorem{theorem}{Theorem}
\newtheorem{corollary}{Corollary}
\newcommand{\MGR}{\mathcal M_{\mathrm{GR}}}
\newcommand{\PhiSrc}{\Phi_{\mathrm{src}}}
\title{V21 P14-T04: Protected Decision on the Status of the Target-Side Congruence}
\author{Gate Chair parent--child synthesis}
\date{2026-07-31}

\begin{document}
\maketitle

\begin{center}
\fbox{\parbox{0.93\linewidth}{\textbf{Protected verdict.}
Within the exact-GR benchmark, the field $u^\mu$ is an
\emph{interpretive representative observer congruence}: admissible,
generally nonunique target-side structure used to decompose a fixed GR
solution.  Current source ontology and dynamics do not identify it with
$\PhiSrc$, dynamically select it, or supply it as an independent physical
field.  This verdict does not classify $u^\mu$ as pure gauge and does not
prove that no future source-side bridge or field extension can exist.}}
\end{center}

\section{Authority and fixed evidence}

Approval \texttt{approval-20260731-002} authorizes exactly one verdict-neutral
P14-T04 decision.  It supplies decision authority, not a missing source law,
bridge, action, equation, mode, stability result, coupling, or observational
constraint.  The review holds fixed the registered foundations and geometry
sources and the finalized P4-T05, P5-T08, and P14-T03 evidence.

The canonical source boundary is decisive.  The source pair consists of a
smooth four-dimensional arena carried as primitive debt and an unresolved
source-order or evolution slot $\PhiSrc$.  No source-to-target bridge is
adopted.  The target-side congruence $u^\mu$ is instead an admissible future
timelike unit observer field on a fixed GR spacetime $(\MGR,g)$.

\section{Type and status separation}

\begin{longtable}{p{0.23\linewidth}p{0.28\linewidth}p{0.39\linewidth}}
\toprule
Candidate reading & P14-T04 status & Evidence boundary \\
\midrule
\endhead
Interpretive representative & \textbf{selected} & The geometry source defines
$u^\mu$ as a generally nonunique observer congruence whose kinematics
re-express a fixed GR solution without adding field equations. \\
Pure gauge & not established & No source-side equivalence group, quotient,
orbit theorem, or transformation law identifying congruence choices as gauge
has been derived. \\
Dynamically selected state & not established & Neither the unresolved
$\PhiSrc$ slot nor the proposal-only P5 candidate supplies a lawful
source-to-congruence selector, positive local geometry input, or uniqueness
and stability theorem. \\
Independent dynamical field & not established & No independent action,
Euler--Lagrange equation, constraints, propagating modes, stability analysis,
matter coupling, or observational bound for $u^\mu$ is part of the exact-GR
benchmark or current source package. \\
\bottomrule
\end{longtable}

The selected status is positive but narrow: $u^\mu$ exists as an admissible
target-side representative used for observer-relative interpretation.  It is
not promoted to a source primitive or additional physical degree of freedom.

\section{Representative-independence theorem}

\begin{definition}[Benchmark base and representative observables]
Let $(M,g,\psi)$ be a fixed solution of the target-side Einstein--matter
system, and let
\[
  \mathcal U(g)=\{u\in\Gamma(TM):g(u,u)=-1,\ u\text{ future directed and
  smooth on its stated domain}\}.
\]
A \emph{base observable} is a functional $\mathcal O[g,\psi;P]$ whose
definition contains no selected member of $\mathcal U(g)$ and whose physical
protocol $P$ is fixed independently of the replaceable congruence.  A
\emph{representative observable} is a functional
$\mathcal K[g,\psi;u]$ whose definition explicitly uses $u$, for example the
projector $h^{(u)}_{\mu\nu}=g_{\mu\nu}+u_\mu u_\nu$, expansion, shear,
vorticity, acceleration, or an observer-measured matter decomposition.
\end{definition}

\begin{theorem}[P14T04-THM-TARGET-BENCHMARK-REPRESENTATIVE-INDEPENDENCE-V1]
Fix $(M,g,\psi)$ and its initial or boundary data, let
$u,v\in\mathcal U(g)$ be two admissible congruence representatives on a
common domain, and hold fixed a physical protocol $P$ defined independently
of $u$ and $v$.  Assume the benchmark action varies only $(g,\psi)$, matter
has no additional coupling to the congruence, and no identification with
$\PhiSrc$ or source-to-congruence bridge is used.  Then:
\begin{enumerate}[label=(\roman*)]
  \item every base observable is representative independent,
  \[
     \mathcal O[g,\psi;u,P]=\mathcal O[g,\psi;v,P]
     =\mathcal O[g,\psi;P];
  \]
  \item replacing $u$ by $v$ changes neither the Einstein--matter equations
  satisfied by $(g,\psi)$ nor the solution pair itself;
  \item representative observables may differ,
  $\mathcal K[g,\psi;u]\neq\mathcal K[g,\psi;v]$, and that difference is an
  observer-relative decomposition of the same fixed target solution, not an
  added target equation or source-side evolution law.
\end{enumerate}
\end{theorem}

\begin{proof}
By definition, a base observable factors through the projection
$(g,\psi,u,P)\mapsto(g,\psi,P)$, so substitution of either $u$ or $v$ leaves
its arguments unchanged, proving (i).  The target field equations in the
fixed benchmark are obtained from the Einstein--matter action and contain no
independently varied congruence field; replacing an auxiliary observer
representative therefore changes neither the equations nor the fixed
solution, proving (ii).  Quantities such as
$h^{(u)}$, $\nabla_\mu u^\mu$, $u^\nu\nabla_\nu u^\mu$, shear, and vorticity
contain the representative explicitly, so they can change under
$u\mapsto v$.  Their change is consequently confined to the representative
decomposition, proving (iii).
\end{proof}

\begin{corollary}[Single-law benchmark boundary]
For the fixed exact-GR benchmark, choosing a static, comoving, ZAMO, or
free-fall congruence may change the observer-relative kinematic dictionary but
does not introduce a second gravitational law or a new vector-sector degree
of freedom.
\end{corollary}

\paragraph{Limit of the theorem.}
The theorem is intentionally conditional on the fixed target solution, a
common domain, and base observables whose physical protocols do not use the
replaceable representative.  If $u$ and $v$ denote different detector or
observer worldlines, the physical protocol has changed rather than merely its
representative label; measured energies, velocities, accelerations,
redshifts, or tidal splits may then differ.  A fluid four-velocity may also be
selected by target-side matter equations and initial data without being
selected by $\PhiSrc$.  The theorem is therefore not a claim that all observer
measurements are numerically equal, a source-side gauge-equivalence theorem,
an ontology-uniqueness theorem, or a no-go theorem against future dynamically
selected or independent-field extensions.

\section{Protected decision}

The exact P14-T04 decision code is
\[
\texttt{INTERPRETIVE\_REPRESENTATIVE\_TARGET\_CONGRUENCE\_ONLY}.
\]
It carries the following controlled meanings:
\begin{enumerate}
  \item $u^\mu$ is target-side benchmark structure and is not $\PhiSrc$.
  \item Its admissible choice is generally nonunique under current sources.
  \item Its role is interpretive and representative; it adds a kinematic
  dictionary rather than an independently varied field equation.
  \item Current source dynamics do not select it.
  \item Pure-gauge and independent-field classifications remain unestablished
  rather than rejected as globally impossible.
\end{enumerate}

\section{Downstream burden update}

All future uses must type the object explicitly as either source-side
$\PhiSrc$ or target-side $u^\mu$.  A stronger status requires new positive
evidence:
\begin{itemize}
  \item \textbf{Gauge:} define the congruence-choice equivalence relation or
  group action, its invariant observables, and a quotient or orbit theorem.
  \item \textbf{Dynamically selected:} derive a target-atlas-free source
  bridge, selection functional or law, selected state, and uniqueness,
  covariance, robustness, and stability results.
  \item \textbf{Independent field:} supply an action, equations, constraints,
  degrees of freedom, mode and stability analysis, matter coupling, and
  observational constraints, followed by separately protected adoption.
  \item \textbf{Public language:} lead with ``interpretive representative
  target congruence,'' state its generally nonunique scope, and block
  preferred-frame or added-vector-sector overreads.
\end{itemize}

\section{Distance-to-GR and explicit nonconclusions}

P14-T04 discharges the protected status-ambiguity burden and adds a
target-benchmark theorem.  It does not reduce a physical source-to-GR burden
or update the Distance-to-GR ledger.  In particular, it does not adopt or
modify canonical ontology, derive $\PhiSrc$, construct a source-to-target
bridge, select a preferred frame, add a physical vector field, derive a
metric, matter coupling, or Einstein equations, promote the benchmark, prove
a completed derivation, publish, push, globally reject the program, or show
future source extensions impossible.

\section{Source materials}

The \AE ther-Flow Research Project. (2026, July 24). \emph{V21 P4-T05:
protected ontology-regime decision} [Internal research-control artifact].

The \AE ther-Flow Research Project. (2026, July 26). \emph{P5 source-dynamics
milestone synthesis} [Internal research-control artifact].

The \AE ther-Flow Research Project. (2026, July 31). \emph{Foundations of the
\AE ther-flow interpretation} and \emph{\AE ther-flow geometry} [Registered
canonical ontology sources].

\end{document}
